% संपूर्ण मराठी ग्रंथातील प्रकरण 58: संचांची क्रमिक संकल्पना
% हा संपादनयोग्य स्रोतखंड आहे. संकलनासाठी संपूर्ण स्रोतसंग्रहातील
% अक्षररूपे, पूर्वभाग, चिन्हव्याख्या आणि आकृती-स्रोत आवश्यक आहेत.
% मूळ: Open Logic Project; मजकूर आणि रूपांतर CC BY 4.0.
% यंत्रानुवाद, दुरुस्ती आणि तपासणी: OpenAI Codex —
% GPT-5.6 Sol आणि GPT-6 Sol, दोन्ही Ultra effort.
% दुरुस्ती-पुनरावलोकन: GPT-6.1 Sol, Ultra effort.
\chapter{संचांची क्रमिक संकल्पना}\label{sth:story:chap}




\section{विस्तारात्मकता}\label{sth:story:extensionality:sec}

सर्वांत आधी सांगायचे म्हणजे
संचांची ओळख त्यांच्या
घटक वरून ठरते.
अधिक नेमकेपणाने:

\begin{axiom}[विस्तारात्मकता]
$A$ आणि~$B$ या संचांत
तेच घटक
असतील, तर~$A$
आणि~$B$ हा
एकच संच आहे.
\[
  \lforall[A][\lforall[B][(\lforall[x][(x \in A \liff x \in B)] \lif
  \eq[A][B])]]
\]
\end{axiom}

\ref{sfr:::part}
या संपूर्ण भागात
आपण हे गृहीत धरले
होते. तरी ते
पुन्हा सांगण्यासारखे
आहे. विस्तारात्मकतेचे
स्वयंसिद्धक संचाची
ओळख त्याच्या
घटक वरून
ठरते ही मूलभूत
कल्पना व्यक्त
करते. (म्हणून
संचांची तुलना
\emph{संकल्पनांशी}
करता येते:
अगदी त्याच
वस्तू अनेक
भिन्न संकल्पनांच्या
अंतर्गत येऊ
शकतात.)

हे तत्त्व का
स्वीकारावे? कारण
विस्तारात्मकता
नाकारणे म्हणजे
\emph{संच उपपत्ती}
म्हणता येईल
अशा गोष्टीचाच
त्याग करणे,
असे म्हणणे
विश्वसनीय वाटते.
संच उपपत्ती
म्हणजे विस्तारात्मक
संग्रहांचीच
उपपत्ती.

पण \ref{sth:::part}
मध्ये खरे आव्हान
\emph{कोणते संच
अस्तित्वात आहेत}
हे सांगणारी
तत्त्वे ठरवणे
आहे. या प्रश्नाचे
एकमेव खरोखर
‘‘उघड’’ वाटणारे
उत्तर सिद्धपणे
चुकीचे ठरते.







\section{रसेलची विरोधापत्ती (पुन्हा)}\label{sth:story:rus:sec}

\ref{sfr:::part}
मध्ये आपण अनौपचारिक
संचसिद्धान्त वापरला.
पण एका \emph{अतिशय}
भोळ्या कल्पनेनुसार
संच म्हणजे विधेयांचे
विस्तारच असतात.
या कल्पनेतून पुढील
तत्त्व स्वीकारावे लागेल:

\begin{defish}
  \emph{अनौपचारिक
  गुणधर्माधारित संचरचना.}
  कोणत्याही सूत्र~$\phi$
  साठी
  $\Setabs{x}{\phi(x)}$
  अस्तित्वात असतो.
\end{defish}

हे तत्त्व मोहक
असले, तरी ते
सिद्धपणे विसंगत
आहे. \ref{sfr:set:rus:sec}
मध्ये आपण हे
पाहिले आहे.
पण निष्कर्ष
इतका महत्त्वाचा
आणि इतका
सरळ आहे
की तो
शब्दशः पुन्हा
मांडण्यासारखा आहे.

\begin{thm}[रसेलची विरोधापत्ती]
$R = \Setabs{x}{x \notin x}$
असा कोणताही संच नाही.
\end{thm}

\begin{proof}
$R = \Setabs{x}{x \notin x}$
अस्तित्वात असेल, तर
$R \in R$ तेव्हा
आणि केवळ तेव्हाच
$R \notin R$;
हा व्याघात आहे.
\end{proof}

रसेल यांनी हा
निष्कर्ष जून १९०१
मध्ये शोधला. (मात्र
त्यांनी विरोधापत्ती
आपण आत्ताच
मांडलेल्या रूपात
दिलेली नव्हती,
कारण ते
\emph{Grundgesetze}
मध्ये मांडलेल्या
फ्रेग यांच्या
संच उपपत्तीचा
विचार करत होते.
\ref{sth:story:blv:sec}
मध्ये आपण
याकडे परत
येऊ.) फ्रेग
यांच्या प्रणालीतील
विसंगती स्पष्ट
करण्यासाठी रसेल
यांनी त्यांना
१६ जून १९०२
रोजी पत्र
लिहिले. त्या
पत्रव्यवहारासाठी
आणि थोड्या
पार्श्वभूमीसाठी
van Heijenoort (1967, पृ.~124--8)
पाहा.

या दोन ओळींची
सिद्धता \emph{शुद्ध
तर्कशास्त्रातून}
मिळते, यावर
भर द्यायला हवा.
मान्य आहे की
$\Setabs{x}{x \notin x}$
या लेखनात
आपण विस्तारात्मकतेचे
एक (तार्किकेतर?)
स्वयंसिद्धक अप्रत्यक्षपणे
वापरले आहे:
$\Setabs{x}{\phi(x)}$
म्हणजे, असा संच
असेल तर,
$\phi$ पूर्ण
करणाऱ्या वस्तूंचा
\emph{एकमेव}
(विस्तारात्मकतेमुळे)
संच. पण निष्कर्ष
पुढीलप्रमाणे
मांडल्यास
विस्तारात्मकतेचा
आभासही टाळता
येतो: \emph{ज्या
संचांचे स्वतःत
सदस्यत्व नाही,
नेमके तेच
ज्याचे सदस्य
आहेत असा
संच नाही}.
याचा संचांशीही
फारसा विशेष
संबंध नाही.
रसेल यांनी
स्वतः नमूद
केल्याप्रमाणे,
अगदी तसाच
युक्तिवाद
दाखवतो:
\emph{जे पुरुष
स्वतःची दाढी
करत नाहीत,
नेमक्या त्याच
पुरुषांची दाढी
करणारा पुरुष
नाही}. किंवा:
\emph{जी पग
जातीची कुत्री
स्वतःचा वास
घेत नाहीत,
नेमक्या त्याच
पग कुत्र्यांचा
वास घेणारा
पग कुत्रा
नाही}. इत्यादी.
सांकेतिक रूपात
हा निष्कर्ष
इतकाच:
\[
\lnot \exists x \forall z(Rzx \liff \lnot R zz).
\]
हे प्रथम-क्रम
तर्कशास्त्राचे
प्रमेय(-रूपबंध)
आहे. त्यामुळे
संच उपपत्तीत
फक्त फेरबदल
करून रसेलची
विरोधापत्ती
टाळता येत
नाही: संच
उपपत्तीपर्यंत
\emph{पोहोचण्याआधीच}
ती उद्भवते.
आपण (अभिजात)
प्रथम-क्रम
तर्कशास्त्र
वापरणार असू,
तर
$R = \Setabs{x}{x\notin x}$
असा संच नाही
हे \emph{स्वीकारावेच}
लागते.

सार असा:
विसंगती
\emph{टाळूनही}
अनौपचारिक
गुणधर्माधारित
संचरचना
स्वीकारायची
असेल, तर
केवळ \emph{संच
उपपत्तीत}
किरकोळ फेरबदल
पुरेसे नाहीत.
त्याऐवजी
\emph{तर्कशास्त्र}
मुळापासून
बदलावे लागेल.

अर्थात,
अभिजातेतर
तर्कशास्त्रे
वापरणाऱ्या
संच उपपत्ती
मांडल्या गेल्या
आहेत. पण
त्या, सौम्य
शब्दांत सांगायचे
तर, मानक
नाहीत.
रसेलच्या
विरोधापत्तीला
सामोरे जाण्याची
मानक पद्धत
म्हणजे तिला
संच अस्तित्वात
नसल्याची
सरळ सिद्धता
मानणे, आणि
मग या निष्कर्षासह
काम कसे करावे
हे शिकणे.
आपण हीच
पद्धत वापरू.









\section{स्वसमावेशरहित आणि स्वसमावेशक व्याख्या}\label{sth:story:predicative:sec}

रसेलचा संच~$R$ हा
$\Setabs{x}{x \notin
x}$ वापरून परिभाषित
केला होता. अधिक
स्पष्टपणे सांगायचे
तर~$R$ म्हणजे
स्वतःचे सदस्य
नसलेले, आणि
फक्त असेच,
सर्व संच सदस्य
म्हणून असलेला
संच ठरला असता.
म्हणून~$R$ ची
व्याख्या करताना
आपण~$R$
(अस्तित्वात असता)
ज्या प्रांतात
आला असता
त्या प्रांतावर
परिमाणन करतो.

ही \emph{स्वसमावेशक}
व्याख्या आहे.
सर्वसाधारणपणे,
परिभाषित होणारी
वस्तू ज्या
प्रांतात समाविष्ट
आहे त्या
प्रांतावर
परिमाणन करणारी
व्याख्या तेव्हा
आणि केवळ तेव्हाच
स्वसमावेशक
असते असे
म्हणता येईल.

या विरोधापत्तींच्या
पार्श्वभूमीवर
व्हाइटहेड, रसेल,
प्वाँकारे आणि
वाइल यांनी
अशा स्वसमावेशक
व्याख्यांना
‘‘दुष्टचक्र निर्माण
करणाऱ्या’’ म्हणून
नाकारले:
\begin{quote}
  टाळायच्या
  विरोधापत्तींचे
  विश्लेषण
  दाखवते की
  त्या सर्व
  एका प्रकारच्या
  दुष्टचक्रातून
  निर्माण होतात.
  एखाद्या
  वस्तुसंग्रहात
  असे सदस्य
  असू शकतात
  जे फक्त
  त्या संपूर्ण
  संग्रहाच्या
  साहाय्यानेच
  परिभाषित
  करता येतात,
  असे मानल्याने
  संबंधित दुष्टचक्रे
  तयार होतात
  [\ldots. \textparagraph]

  अनुचित समग्रता
  टाळू देणारे
  तत्त्व असे
  मांडता येते:
  ‘संग्रहातील
  \emph{सर्वांचा}
  उल्लेख ज्यात
  येतो, ती
  वस्तू त्या
  संग्रहातील
  एक वस्तू
  असू नये’;
  किंवा उलट:
  ‘एखाद्या
  संग्रहाला
  समग्रता असती
  तर त्यातील
  काही सदस्य
  फक्त त्या
  समग्रतेच्या
  साहाय्यानेच
  परिभाषित
  करता आले
  असते, तर
  त्या संग्रहाची
  समग्रता
  अस्तित्वात
  नाही.’ याला
  आपण
  ‘दुष्टचक्र
  तत्त्व’
  म्हणू,
  कारण
  अनुचित
  समग्रता
  गृहीत
  धरल्याने
  उद्भवणारी
  दुष्टचक्रे
  ते टाळू
  देते.
  (Whitehead आणि Russell, 1910, पृ.~37)
\end{quote}
त्यांच्याप्रमाणे
\emph{दुष्टचक्र तत्त्व}
स्वीकारले, तर
\ref{sth:story:rus:sec}
मधील विनाशकारी
अनौपचारिक
गुणधर्माधारित
संचरचना-रूपबंधाऐवजी
पुढील प्रकारचे
तत्त्व वापरण्याचा
प्रयत्न करता
येईल:

\begin{defish}
\emph{स्वसमावेशरहित
गुणधर्माधारित
संचरचना.}
फक्त आधीच्या स्तरातील संचांवर परिमाणन करणाऱ्या
प्रत्येक सूत्र~$\phi$ साठी, त्याच स्तरातील $x$ घेऊन बनणारा
संच$^\prime$ $\Setabs{x}{\phi(x)}$ पुढील स्तरात अस्तित्वात असतो.
त्याचे सदस्य आधीच्या स्तरातील संचच असतात.
ही स्तरांची अनौपचारिक रूपरेषा आहे;
पूर्ण औपचारिक भाषा आणि स्वयंसिद्धे येथे दिलेली नाहीत.
\end{defish}

संच$^{\prime}$ यांची निर्मिती आधीच्या स्तरातील संचांपासून
केल्यास नव्याने परिभाषित होणारी वस्तू त्या सदस्य-प्रांतात येत नाही.
त्यामुळे वरचा रसेलचा स्वसदस्यत्वाचा युक्तिवाद तसाच लागू होत नाही.
एवढ्यावरून कोणत्याही संपूर्ण औपचारिक उपपत्तीची
सुसंगतता सिद्ध होत नाही; त्यासाठी तिची अचूक भाषा,
स्वयंसिद्धे आणि स्वतंत्र सुसंगतता-युक्तिवाद आवश्यक आहेत.

दुर्दैवाने,
स्वसमावेशरहित
गुणधर्माधारित
संचरचना फारशी
\emph{व्यापक}
नाही. कारण
तिच्यामुळे
संच$^\prime$
या नव्या
वस्तू मिळतात.
म्हणून संच$^\prime$
यांच्यावर
परिमाणन करणारी
सूत्रेही विचारात
घ्यावी लागतील.
त्या सूत्रांसाठी सदस्य-प्रांत आणि नवीन वस्तूंचा स्तर
पुन्हा एकच केला आणि त्या स्तरात अमर्याद गुणधर्माधारित
संचरचना मानली, तर स्वतःचे
सदस्य नसलेल्या
सर्व संच$^\prime$
यांचा संच$^\prime$
पाहताच रसेलची
विरोधापत्ती
पुन्हा उद्भवेल.
त्यामुळे याच
विचारानुसार
संच$^\prime$
यांच्यावर
परिमाणन करणारे
सूत्र अनुरूप
संच$^{\prime\prime}$
देते असे
म्हणावे लागेल.
पुढे
संच$^{\prime\prime\prime}$,
संच$^{\prime\prime\prime\prime}$
इत्यादी लागतील.
अशी कितीतरी
प्राइम चिन्हे
टाळण्यासाठी
यांना संच$_0$,
संच$_1$,
संच$_2$,
संच$_3$,
संच$_4$,\ldots
असे मानणे
सोपे ठरेल.
यातून (सरळ)
प्रकार-उपपत्तीचा
मार्ग मिळतो.

अशा उपपत्तीवर
काही उघड
आक्षेप आहेत
(मात्र ते
\emph{निर्णायक}
आहेतच असे
स्पष्ट नाही).
थोडक्यात:
तिचा वापर
जड जातो;
ती अनेक
भिन्न प्रकारच्या
वस्तू गृहीत
धरते; आणि
शेवटी
स्वसमावेशक
व्याख्या
\emph{खरोखर
इतक्या वाईट}
आहेत का,
हेही स्पष्ट
नाही.

शेवटचा मुद्दा
स्पष्ट करण्यासाठी
रॅम्झी
यांची पुढील
टिप्पणी पाहा:
\begin{quote}
  एखाद्या
  गटातील
  सर्वांत
  उंच माणूस
  म्हणून आपण
  एखाद्याचा
  उल्लेख करू
  शकतो. असे
  केल्याने तो
  स्वतः ज्या
  समग्रतेचा
  सदस्य आहे
  तिच्याच
  साहाय्याने
  त्याची ओळख
  पटते; तरीही
  त्यात दुष्टचक्र
  नसते.
  (Ramsey, 1925)
\end{quote}
रॅम्झी यांचा
मुद्दा असा
की ‘‘गटातील
सर्वांत उंच
माणूस’’ ही
\emph{स्वसमावेशक}
व्याख्या आहे;
पण ती
स्पष्टपणे
निर्दोष आहे.

त्यावर असे
उत्तर देता
येईल की
या उदाहरणात
सर्वांत उंच
व्यक्तीला
\emph{स्वसमावेशरहित}
मार्गानेही
ओळखता येईल.
उदाहरणार्थ,
आपण त्या
माणसाकडे
सरळ बोट
दाखवू
शकतो. मग
स्वसमावेशक
व्याख्यांवरील
आक्षेप फक्त
\emph{स्वसमावेशक
मार्गानेच}
ओळखता येणाऱ्या
वस्तूंपुरता
मर्यादित करावा
लागेल. पण
अशा
‘‘मूलभूत
स्वसमावेशकतेत’’
नेमके वाईट
काय आहे,
हेही पुढे
स्पष्ट करावे
लागेल.\footnote{
  अधिक माहितीसाठी
  Linnebo (2010)
  पाहा.}

वस्तू
\emph{निर्माण}
करण्यासाठी
आपल्या
व्याख्यांनी
कृतीसूचना
द्यावी असे
अपेक्षित
असेल, तर
स्वसमावेशक
व्याख्या
निश्चितच
मोठी अडचण
ठरतात.
कारण अशा
व्याख्येवरून
वस्तू बनवताना
खरेच
दुष्टचक्रात
अडकावे
लागेल:
स्वसमावेशकपणे
निर्दिष्ट
वस्तू निर्माण
करण्यासाठी
\emph{आधी} सर्व
वस्तू
(त्यात ती
निर्दिष्ट
वस्तूही येते)
निर्माण कराव्या
लागतील,
कारण तिची
व्याख्या
सर्व वस्तूंच्या
संदर्भात आहे.
म्हणजे त्या
वस्तूची
निर्मिती
होण्यापूर्वीच
तिची निर्मिती
करावी लागेल.
पण पुन्हा,
‘‘मूलभूत
स्वसमावेशक’’
रीतीने
निर्दिष्ट
संचांवरील
हा गंभीर
आक्षेप
तेव्हाच
ठरतो,
जेव्हा
संच म्हणजे
आपण
\emph{निर्माण}
करत असलेल्या
वस्तू आहेत
असे मानतो.
आपण
(बहुधा)
तसे मानत
नाही.

त्यामुळे,
चांगल्या किंवा
वाईट कारणांनी,
रूढ झालेल्या
पद्धतीत
(अ)स्वसमावेशकतेबाबत
कठोर भूमिका
घेतलेली नाही.
त्याऐवजी
आता संचयी-क्रमिक
पद्धत मानली
जाते. शेवटी
तिच्यामुळे
संचांचे
‘‘टप्प्यां’’मध्ये
स्तरीकरण करता
येते---स्वसमावेशरहित
पद्धतीत वस्तू
संच$_0$,
संच$_1$,
संच$_2$,
\ldots मध्ये
स्तरीकृत
होतात तसेच
\emph{थोड्याफार}
प्रमाणात.
मात्र या
टप्प्यांतील
संचांच्या
मूलभूत प्रकारांत
कोणताही
फरक गृहीत
धरला जात
नाही.








\section{संचयी-क्रमिक पद्धत}\label{sth:story:approach:sec}

संचांचे टप्प्यांमध्ये
स्तरीकरण कसे
करायचे, याचे
थोडे अधिक
तपशीलवार
वर्णन असे:
\begin{quote}
  संच \emph{टप्प्यांमध्ये}
  तयार होतात.
  प्रत्येक टप्पा~$S$
  यासाठी, $S$ च्या \emph{आधी}
  कोणते टप्पे येतात, हे निश्चित असते.
  $S$ टप्प्यावर,
  $S$ पूर्वीच्या
  टप्प्यांवर
  तयार झालेले
  संच घेऊन
  बनलेल्या
  प्रत्येक
  संग्रहाचा
  संच तयार
  होतो.
  टप्प्यांवर
  तयार झालेले
  संच सोडून
  इतर संच
  नाहीत.
  (Shoenfield, 1977, पृ.~323)
\end{quote}
ही \emph{संचांची
संचयी-क्रमिक
संकल्पना}
याची रूपरेषा
आहे. \ref{sth:::part}
मध्ये मांडल्या
जाणाऱ्या
औपचारिक
संच उपपत्तीला
ती आधार
देईल.

थोडा अधिक
तपशील पाहू.
शोनफिल्ड
यांच्या
वर्णनानुसार,
प्रत्येक
टप्प्यावर
मागील
टप्प्यांपासून
आपल्याला
उपलब्ध
असलेल्या
संचांपासून
नवे संच
तयार होतात.
म्हणून
आरंभीच्या
टप्प्यावर,
म्हणजे
टप्पा~$0$
येथे,
\emph{आधीचा}
टप्पाच
नसतो;
\emph{मग तर}
पूर्वीच्या
टप्प्यांतून
आपल्याला
संचही
उपलब्ध
नसतात.\footnote{
  पहिला
  टप्पा
  \emph{असतो}
  असे का
  गृहीत
  धरावे?
  \ref{sth:z:story:sec}
  मधील
  \stagesord{}
  जवळची
  तळटीप
  पाहा.}
त्यामुळे
आपण
एकच संच
तयार करतो:
एकही
घटक
नसलेला
रिक्त संच
$\emptyset$.
टप्पा~$1$
येथे
मागील
टप्प्यांपासून
बरोबर
एक संच
उपलब्ध
असतो,
म्हणून
फक्त एक
नवा संच
$\{\emptyset\}$
तयार होतो.
टप्पा~$2$
येथे मागील
टप्प्यांपासून
दोन संच
उपलब्ध
असतात;
आपण दोन
नवे संच
$\{\{\emptyset\}\}$
आणि
$\{\emptyset,
\{\emptyset\}\}$
तयार करतो.
टप्पा~$3$
येथे मागील
टप्प्यांपासून
चार संच
उपलब्ध
असतात;
म्हणून
बारा
नवे संच
तयार
होतात\ldots.
त्यामुळे
संचांचे
संचयी-क्रमिक
चित्र
साधारण
असे
दिसेल
(अंक टप्पे
दर्शवतात):
\begin{center}
  \begin{tikzpicture}[scale=0.6]
  \tikzset{cut_here/.style={densely dotted}}
  \draw (-4,6) -- (-1, 0)-- (2,6);
  \draw[cut_here] (-5,8)--(-4,6);
  \draw[cut_here] (2,6)--(3,8);
  \draw(-1.5, 1)--(-0.5, 1);
  \draw(-2, 2)--(0, 2);
  \draw(-2.5, 3)--(0.5,3);
  \draw(-3, 4)--(1, 4);
  \draw(-3.5, 5)--(1.5, 5);
  \draw(-4, 6)--(2, 6);
  \node[label] at (0, 0) {\small 0};
  \node[label] at (0.5, 1) {\small 1};
  \node[label] at (1, 2) {\small 2};
  \node[label] at (1.5, 3) {\small 3};
  \node[label] at (2, 4) {\small 4};
  \node[label] at (2.5, 5) {\small 5};
  \node[label] at (3, 6) {\small 6};
  \end{tikzpicture}
\end{center}
मग ही मांडणी
का स्वीकारावी?

एक कारण:
ती चांगली
आणि हाताळायला
सुलभ मांडणी
आहे. सर्वांत
उघड वाटणारी
मांडणी,
म्हणजे
अनौपचारिक
गुणधर्माधारित
संचरचना,
अपयशी ठरल्यावर
आपल्याला
एखादी चांगली
मांडणी हवीच.

पण ही मांडणी \emph{फक्त} सुलभ नाही.
ती रसेलच्या विरोधापत्तीला कारणीभूत असलेली
अमर्याद संचरचना का मानू नये हे स्पष्ट करते
आणि संच उपपत्तीची स्वयंसिद्धे निवडण्यास आधार देते.
खालील युक्तिवाद हा रसेलचा संच का बनत नाही यासाठी आहे;
या मांडणीवर आधारित कोणतीही उपपत्ती \emph{सुसंगत} असते,
असे तो सिद्ध करत नाही. पूर्ण सुसंगततेसाठी
नेमकी स्वयंसिद्धे आणि स्वतंत्र युक्तिवाद आवश्यक असतो.

संचांची
संचयी-क्रमिक
संकल्पना
मानल्यास
संच
टप्प्याटप्प्याने
बनतात;
त्यांचे
घटक
म्हणजे
\emph{आधीच}
उपलब्ध
असलेल्या
वस्तू असतात.
म्हणून
कोणत्याही
टप्प्या~$S$
साठी
पुढील संच
तयार करता
येतो:
\[
  R_S = \Setabs{x}{x \notin x \text{ आणि $x$ हा $S$ पूर्वी उपलब्ध होता}}
\]
रसेलची
विरोधापत्ती
सिद्ध करताना
वापरलेला
युक्तिवाद
दाखवतो
की~$R_S$
हा संच
स्वतः
टप्पा~$S$
पूर्वी
उपलब्ध
नसतो.
आणि
त्यात
व्याघात
नाही.
शिवाय
संचांची
संचयी-क्रमिक
संकल्पना
स्वीकारल्यास
रसेलचा
मूळ संच
तयार होईल
अशी
\emph{अपेक्षाच}
ठेवायला
नको होती.
कारण तो
‘‘कधीही
उपलब्ध
होतील’’
अशा
स्वतःचे
सदस्य
नसलेल्या
सर्व
संचांचा
संच
असता.
थोडक्यात,
संचयी-क्रमिक
मांडणी
मानल्यावर
रसेलचा
संच तयार
करता
येत
नाही
ही
(सिद्ध
होणारी)
गोष्ट
\emph{आश्चर्यकारक}
नाही;
याचाच
\emph{अंदाज}
आपण
केला
असता.













\section{मूलघटक असावेत का?}\label{sth:story:urelements:sec}

पुढील काही प्रकरणांत आपण संचांच्या संचयी-क्रमिक संकल्पनेतून स्वयंसिद्धके मिळवण्याचा प्रयत्न करू. पण पुढे जाण्याआधी \emph{मूलघटकांबद्दल} आणखी काही सांगणे गरजेचे आहे.

\ref{sth:story:approach:sec} मधील चित्रात आपण जुन्या \emph{संचांपासून} नवे संचच तयार करत होतो. पण काही \emph{संच नसलेल्या} वस्तू---गायी, डुकरे, वाळूचे कण किंवा इतर काही---संचांचे घटक असू शकतात, असे मानायचे असेल. तसे असल्यास सुरुवातीला काही मूलभूत वस्तू, म्हणजे \emph{मूलघटक}, उपलब्ध मानू. मग प्रत्येक टप्प्या~$S$ वर त्या मूलघटकांपासून किंवा $S$ पूर्वीच्या टप्प्यांवर तयार झालेल्या संचांपासून बनणारे ``सर्व शक्य'' संच तयार होतात, असे म्हणू; यात दोन्ही प्रकारच्या वस्तू एकत्रही असू शकतात. अशा वेळी चित्र साधारण असे दिसेल:
\begin{center}
	\begin{tikzpicture}[scale=0.6]
	\tikzset{cut_here/.style={densely dotted}}
	\draw (-4,6) -- (-1, 0) -- (1,0) -- (4,6);
	\draw[cut_here] (-5,8)--(-4,6);
	\draw[cut_here] (5,8)--(4,6);
	\draw(-1.5, 1)--(1.5, 1);
	\draw(-2, 2)--(2, 2);
	\draw(-2.5, 3)--(2.5,3);
	\draw(-3, 4)--(3, 4);
	\draw(-3.5, 5)--(3.5, 5);
	\draw(-4, 6)--(4, 6);
	\node[label] at (2, 0) {\small 0};
	\node[label] at (2.5, 1) {\small 1};
	\node[label] at (3, 2) {\small 2};
	\node[label] at (3.5, 3) {\small 3};
	\node[label] at (4, 4) {\small 4};
	\node[label] at (4.5, 5) {\small 5};
	\node[label] at (5, 6) {\small 6};
	\end{tikzpicture}
\end{center}
आता आपल्याला एक निर्णय घ्यायचा आहे: \emph{मूलघटकांना परवानगी द्यावी का?}

तात्त्विक दृष्टीने आपल्या सिद्धान्तात मूलघटक समाविष्ट करणे सयुक्तिक आहे. त्यामागचे मुख्य कारण म्हणजे संच उपपत्ती \emph{लागू करता येण्याजोगी} बनवणे. हे स्पष्ट करण्यासाठी \ref{sfr:siz::chap} आठवा: दोन संच $A$ आणि~$B$ यांचा आकार समान, म्हणजे $\cardeq{A}{B}$, तेव्हा आणि केवळ तेव्हाच त्यांच्यात एकास-एक व आच्छादक फलन असते. आता शेतातल्या गायी आणि गोठ्यातली डुकरे हे दोन्ही संच बनवत असतील, तर ``गायी जितक्या आहेत तितकीच डुकरे आहेत'' या विधानाचा संच उपपत्तीच्या आधारे विचार करता येतो. पण मूलघटकांना मनाई केली आणि गायी व डुकरे यांचे संच \emph{बनतच नाहीत} असे मानले, तर हा विचार करता येणार नाही. खरे तर संचांशिवाय इतर कशालाही संच उपपत्ती लागू करण्याचा सरळ मार्ग आपल्याकडे उरणार नाही. (मूलघटक समाविष्ट करण्याची आणखी कारणे Potter (2004, पृ.~vi, 24, 50--1) येथे पाहा.)

मात्र गणितात मूलघटकांना परवानगी देणे फारसे प्रचलित नाही. त्याचे एक कारण म्हणजे मूलघटकांशिवाय संच उपपत्ती मांडणे \emph{अगदी किंचित} सोपे असते. पण कधीकधी त्यांना वगळण्याचे अधिक लक्षवेधी समर्थनही आढळते:
\begin{quote}
  संच उपपत्ती हा गणिताचा पाया आहे या मतानुसार, केवळ संचांविषयी बोलून आपल्याला संपूर्ण गणित व्यक्त करता आले पाहिजे. म्हणून आपल्या चलांची व्याप्ती गायी आणि डुकरे यांसारख्या वस्तूंवर असता कामा नये.

  (Kunen, 1980, पृ.~8)
\end{quote}
म्हणून उपयोज्यतेला महत्त्व दिल्यास मूलघटक \emph{समाविष्ट} करावेसे वाटतील; पण गणित शुद्ध संच उपपत्तीत आणण्याच्या पायाभूत ध्येयाला महत्त्व दिल्यास त्यांना \emph{वगळावेसे} वाटेल. थोडा आळसही मूलघटक वगळण्याकडे झुकवतो.

आपण सर्वांत सोपा मार्ग निवडू. पण त्यामागे काही प्रमाणात शैक्षणिक कारणही आहे. संच उपपत्तीच्या तत्त्वज्ञानाचा अभ्यास सुरू करण्यासाठी तुम्हाला तिच्या ज्या मूलभूत गोष्टी लागतील, त्यांची ओळख करून देणे हे आपले उद्दिष्ट आहे. या विषयावरील बहुतेक तात्त्विक साहित्य मूलघटक \emph{नसलेल्या} संच उपपत्तीची चर्चा करते. म्हणून त्या साहित्याशी साधारण जुळून राहिल्यास हे पुस्तक अधिक उपयोगी ठरेल.









\section{परिशिष्ट: फ्रेगे यांचा मूलनियम पाच}\label{sth:story:blv:sec}

\ref{sth:story:rus:sec} मध्ये आपण पाहिले की रसेल यांनी आपली विरोधापत्ती फ्रेगे यांनी \emph{Grundgesetze} मध्ये मांडलेल्या प्रणालीसमोरील अडचण म्हणून सांगितली. फ्रेगे यांच्या प्रणालीत अनौपचारिक गुणधर्माधारित संचरचनेची थेट मांडणी नव्हती. म्हणून या परिशिष्टात त्या प्रणालीत \emph{नेमके काय} होते, त्याचा अनौपचारिक गुणधर्माधारित संचरचनेशी कसा संबंध आहे आणि रसेल यांच्या विरोधापत्तीशी कसा संबंध आहे, हे अगदी थोडक्यात पाहू.

फ्रेगे यांची प्रणाली \emph{द्वितीय-क्रमाची} आहे. तिचा हेतू \emph{संकल्पनेचा विस्तार} ही संकल्पना औपचारिकरीत्या मांडणे हा होता.\footnote{काटेकोरपणे सांगायचे तर फ्रेगे यांनी अधिक व्यापक संकल्पना, म्हणजे फलनाची ``मूल्यव्याप्ती'', औपचारिकरीत्या मांडण्याचा प्रयत्न केला. संकल्पनांचे विस्तार हे त्या व्यापक संकल्पनेचे विशेष उदाहरण आहे. तपशीलांसाठी Heck (2012, पृ.~8--9) पाहा.} फ्रेगे यांच्या संकेतपद्धतीच्या धर्तीवर $\fregeext{x}{F(x)}$ हे \emph{संकल्पना~$F$ चा विस्तार} दर्शवेल. हे साधन ``$F$'' या \emph{विधेयाचे} (प्रथम-क्रमातील) ``$\fregeext{x}{F(x)}$'' या \emph{पदात} रूपांतर करते. या साधनाच्या आधारे फ्रेगे यांनी सदस्यत्वाची पुढील \emph{व्याख्या} दिली:
\[
	a \in b =_\text{df} \exists G(b = \fregeext{x}{G(x)} \land Ga)
\]
स्थूलमानाने, $a \in b$ तेव्हा आणि केवळ तेव्हाच जेव्हा $a$ अशी एखादी संकल्पना पूर्ण करते की तिचा विस्तार $b$ आहे. (येथे ``$\exists G$'' हा संख्यापक द्वितीय-क्रमाचा आहे, हे लक्षात घ्या.) फ्रेगे यांनी \emph{मूलनियम पाच} म्हणून ओळखले जाणारे पुढील तत्त्वही मानले:
$$\fregeext{x}{F(x)} = \fregeext{x}{G(x)} \liff \forall x (Fx \liff Gx)$$
स्थूलमानाने, दोन संकल्पनांचे विस्तार एकच असतात तेव्हा आणि केवळ तेव्हाच त्या नेमक्या त्याच वस्तूंना लागू पडतात. (पुन्हा, ``$F$'' आणि ``$G$'' ही दोन्ही विधेयस्थानी आहेत.) आता एक साधे तत्त्व सदस्यत्वाचा संबंध गुणधर्मपूर्तीशी जोडते:

\begin{lem}[\emph{Grundgesetze} मध्ये]\label{sth:story:blv:lem:Fregeextensions}
$\forall F \forall a(a \in \fregeext{x}{F(x)} \liff Fa)$
\end{lem}

\begin{proof}
$F$ आणि $a$ निश्चित करू. सदस्यत्वाच्या व्याख्येनुसार $a \in \fregeext{x}{F(x)}$ तेव्हा आणि केवळ तेव्हाच जेव्हा $\exists G(\fregeext{x}{F(x)}
= \fregeext{x}{G(x)} \land Ga)$; मूलनियम पाचनुसार तेव्हा आणि केवळ तेव्हाच जेव्हा $\exists G(\forall x(Fx \liff Gx) \land Ga)$; आणि प्राथमिक द्वितीय-क्रम तर्कशास्त्रानुसार तेव्हा आणि केवळ तेव्हाच जेव्हा $Fa$.
\end{proof}

यातून अनौपचारिक गुणधर्माधारित संचरचना जवळजवळ लगेच मिळते:

\begin{lem}[\emph{Grundgesetze} मध्ये.]
$\forall F \exists s \forall a (a \in s \liff Fa)$
\end{lem}

\begin{proof}
$F$ निश्चित करू. मग \ref{sth:story:blv:lem:Fregeextensions} वरून $\forall a (a \in
\fregeext{x}{F(x)} \liff Fa)$ मिळते. म्हणून अस्तित्ववाचक सामान्यीकरणाने $\exists s\forall a(a \in s \liff Fa)$ मिळते. $F$ स्वैर असल्यामुळे अपेक्षित निष्कर्ष मिळतो.
\end{proof}

येथे पूर्ण द्वितीय-क्रम गुणधर्माधारित संकल्पनारचना
वापरून $\exists F\forall x(Fx\liff x\notin x)$ हे गृहीतक मिळते.
सदस्यत्वाच्या वरील व्याख्येत द्वितीय-क्रम संख्यापक आहे;
म्हणून हा संकल्पनारचनेचा उपयोग फक्त प्रथम-क्रम सूत्रांपुरता मर्यादित नाही.
अशी $F$ घेऊन $R=\fregeext{x}{F(x)}$ ठेवले, तर
वरील उपप्रमेयानुसार $R\in R\liff FR$, तर $F$ च्या व्याख्येनुसार
$FR\liff R\notin R$. त्यामुळे रसेलची विरोधापत्ती मिळते.
या युक्तिवादाला मूलनियम पाचासोबत किमान
वरील विशिष्ट संकल्पनारचनेचे उदाहरण आवश्यक आहे;
पूर्ण रूपबंधामुळे ते उपलब्ध होते.
संकल्पनारचना मर्यादित करून सामान्य द्वितीय-क्रम प्रतिमाने वापरल्यास
मूलनियम पाचाच्या सुसंगततेचा वेगळा प्रश्न उभा राहतो;
प्रत्येक मर्यादित रूपबंध असुसंगत असल्याचा हा युक्तिवाद नाही.



